% !TEX root = main.tex
%% ===========================================================================
%%  本文
%%
%%  章が増えて 1 ファイルが扱いにくくなったら、章ごとにファイルを分けて
%%  main.tex から \include{Introduction} のように読むこともできます。
%%  最初のうちは 1 ファイルのほうが全体を見渡しやすいはずです。
%% ===========================================================================

\chapter{序論}
\label{chap:introduction}

%% 序論の役割は「なぜこの研究をやる必要があったのか」を読者に納得させることです。
%% 次の順で書くと形になります。
%%   1. 大きな背景（太陽・恒星の何が面白いか）
%%   2. これまでに分かっていること（先行研究のレビュー、必ず引用をつける）
%%   3. まだ分かっていないこと（先行研究の限界。ここが一番大事）
%%   4. 本研究の目的（3 を受けて、何をするのか）
%%   5. 本論文の構成（第 2 章では…、第 3 章では…）
%% 「面白いから」だけでは読む人が納得しにくいので、
%% 「ここが分かっていないから」まで書けると強いです。

\section{太陽対流と磁場}
\label{sec:intro_convection}

太陽の内部は、中心から約 $0.7\,R_\odot$ より外側が対流層になっている。
対流はエネルギーを外向きに運ぶと同時に、差動回転と子午面循環を駆動し、
それらがダイナモ機構を通じて 11 年周期の磁場を生み出している
\citep{charbonneau2020}。

\section{先行研究}
\label{sec:intro_previous}

%% レビューは「誰が何をやったか」の羅列ではなく、
%% 「この分野の理解がどう進んできたか」の物語として書きます。
%% 引用は \citep{} と \citet{} を使い分けます。
%%   \citep{hotta2021}  ->（Hotta et al. 2021）  文の最後に添えるとき
%%   \citet{hotta2021}  -> Hotta et al.（2021）  主語にするとき
高解像度の大局的計算が可能になったことで、対流層の乱流構造が詳しく調べられてきた。
\citet{hotta2021} は小スケールダイナモが対流の有効粘性を下げ、
太陽に近い差動回転が再現されることを示した。

\section{本研究の目的}
\label{sec:intro_purpose}

\todo{先行研究で残された問題と、それに対して本研究が何をするのかを書く}

\section{本論文の構成}
\label{sec:intro_outline}

第\ref{chap:method}章では本研究で用いた手法を述べる。
第\ref{chap:results}章で結果を示し、第\ref{chap:discussion}章で議論する。
第\ref{chap:conclusion}章は結論である。


\chapter{手法}
\label{chap:method}

%% 「他人が読んで同じことを再現できるか」が唯一の判断基準です。
%% 数値計算なら：基礎方程式、境界条件、初期条件、解像度、計算領域、
%%              使ったコードとそのバージョン、計算機環境。
%% 観測なら    ：装置、観測日時、露出、較正の手順、使ったソフトのバージョン。
%% 「標準的な手法を用いた」で済ませないでください。何が標準かは 10 年で変わります。

\section{基礎方程式}
\label{sec:method_equations}

本研究では磁気流体力学（MHD）方程式を解く。
連続の式、運動方程式、誘導方程式はそれぞれ
%% 番号を後で参照する式は align、参照しない式は align* を使います。
\begin{align}
  \pdif{\rho}{t} &= -\divergence(\rho\vect{v}), \label{eq:continuity}\\
  \pdif{(\rho\vect{v})}{t}
    &= -\divergence(\rho\vect{v}\vect{v})
       - \grad p + \rho\vect{g}
       + \frac{1}{4\pi}(\curl\vect{B})\times\vect{B}, \label{eq:motion}\\
  \pdif{\vect{B}}{t} &= \curl(\vect{v}\times\vect{B}) \label{eq:induction}
\end{align}
である。ここで $\rho$ は密度、$\vect{v}$ は速度、$p$ は圧力、
$\vect{B}$ は磁場、$\vect{g}$ は重力加速度である。
式\eqref{eq:induction}では磁気拡散を無視した。

\section{計算設定}
\label{sec:method_setup}

\todo{計算領域、解像度、境界条件、初期条件を表にまとめる}

%% 表は booktabs を使うと縦罫線なしの読みやすい形になります。
%% 表のキャプションは表の「上」に置きます。
\begin{table}[tb]
  \centering
  \caption{数値計算の設定。}
  \label{tab:setup}
  \begin{tabular}{lll}
    \toprule
    項目 & 値 & 備考 \\
    \midrule
    計算領域   & $0.71\,R_\odot \le r \le 0.96\,R_\odot$ & 対流層 \\
    格子数     & $512\times1024\times2048$               & $(r,\theta,\phi)$ \\
    回転角速度 & \qty{2.6e-6}{\per\second}               & 太陽の値 \\
    \bottomrule
  \end{tabular}
\end{table}


\chapter{結果}
\label{chap:results}

%% 結果の章では「何が見えたか」だけを書きます。解釈は次の章です。
%% ただし機械的に図の説明を並べるのではなく、
%% 読者が図を見る順番を自分が案内するつもりで書いてください。
%% 図は必ず本文から \ref{} で参照します。本文から触れていない図があったら、
%% その図は要らないか、本文に書き忘れているかのどちらかです。

\section{対流構造}
\label{sec:results_convection}

\todo{得られた結果を書く}

%% 図のキャプションは図の「下」に置きます。
\begin{figure}[tb]
  \centering
  \includegraphics[width=0.7\linewidth]{sample.pdf}
  \caption{図の説明。キャプションだけを読んで何の図か分かるように書くこと。
    軸の量と単位、線や色が何を表すか、どの計算・観測のものかを含める。}
  \label{fig:sample}
\end{figure}

図\ref{fig:sample}に\todo{何を示した図か}を示す。

%% 図を横に並べて (a) (b) と番号を振る例。subcaption パッケージを使う。
\begin{figure}[tb]
  \centering
  \begin{subfigure}{0.48\linewidth}
    \includegraphics[width=\linewidth]{sample.pdf}
    \caption{低解像度}
    \label{fig:comp_low}
  \end{subfigure}\hfill
  \begin{subfigure}{0.48\linewidth}
    \includegraphics[width=\linewidth]{sample.pdf}
    \caption{高解像度}
    \label{fig:comp_high}
  \end{subfigure}
  \caption{解像度による対流構造の違い。個々の図は
    図\ref{fig:comp_low}、図\ref{fig:comp_high} のように参照できる。}
  \label{fig:comp}
\end{figure}


\chapter{議論}
\label{chap:discussion}

%% 議論の章でやることは 3 つです。
%%   1. 結果の解釈：なぜそうなったのかを物理で説明する
%%   2. 先行研究との比較：一致するのか、しないのか。しないならなぜか
%%   3. 限界：この研究で言えないことは何か
%% 3 を書くのを怖がらないでください。限界を自分で把握できていることは
%% 減点ではなく加点になります。

\section{結果の解釈}
\label{sec:discussion_interpretation}

\todo{書く}

\section{先行研究との比較}
\label{sec:discussion_comparison}

\todo{書く}

\section{本研究の限界}
\label{sec:discussion_limitation}

\todo{書く}


\chapter{結論}
\label{chap:conclusion}

%% 結論は「まとめ」ではなく、序論で立てた問いに対する答えを短く言い切る場所です。
%% 新しい図やデータはここでは出しません。
%% 最後に今後の展望を 1 段落だけ書きます。

本研究では\todo{何をしたか}を行った。その結果、以下のことが分かった。

\begin{enumerate}
  \item \todo{結論 1}
  \item \todo{結論 2}
\end{enumerate}

今後は\todo{展望}が課題である。
